\section*{The Propers: Notable and Quotable}

{\small
\begin{itemize}[itemsep=0em,topsep=0.15em,parsep=0pt]
\item \textbf{“God so loved the world” --- transformed into “Austin 3:16.”\enspace\elementref{Gospel}} After defeating Bible-quoting Jake Roberts at the 1996 King of the Ring, “Stone Cold” Steve Austin mocked John 3:16 with “Austin 3:16,” an improvised boast that became a catchphrase, arena sign, blockbuster shirt, and marker of WWE's Attitude Era. A saving verse becomes aggressive anti-verse and pop-cultural brand.

\item \textbf{“God so loved the world” --- sealed beneath Cleopatra's Needle for readers after empire.\enspace\elementref{Gospel}} In 1878 a copy of John~3:16 in 215 languages entered the obelisk's foundation deposit beside imperial coins, a railway timetable, newspapers, weights, and other Victorian objects. The Gospel's “world” became a multilingual message to an unknown future civilization; under an ancient monument, it also quietly outlasts the empire that buried it.

\item \textbf{“A stiff-necked people” --- even a dinner party can become one.\enspace\elementref{First Reading}} The biblical bodily image has entered ordinary English as “stiff-necked”: haughty or stubborn, but also formal or stilted. Merriam-Webster's secondary sense lets an accusation of covenantal obstinacy describe an awkwardly ceremonious social occasion.
\end{itemize}
}

\Needspace{10\baselineskip}
\section*{The Propers: Interpretive Possibilities}

\noindent\emph{The following exploratory proposals were developed in an AI-assisted editorial process. They are attributed to none of the cited authorities and claim no historical compositional intent. Each states its anchors and limits.}

\Needspace{10\baselineskip}
\noindent\textbf{Can presence move from “among us” to “within us” without leaving Sinai behind?\enspace\elementref{First Reading, Collect, Gospel, Communion}}
Moses asks the Lord to travel in the people's midst; the Collect names the Father's sending of Word and sanctifying Spirit; John proclaims the Son sent into the world; and the Communion formula places the Spirit of the Son within human hearts. The connection rests on a change in the grammar of presence---\emph{in the midst}, \emph{into the world}, \emph{into hearts}---not on one text predicting the next. Reading each unit alone can miss this controlled intensification of divine accompaniment: mercy does not merely overlook stiff-neckedness but creates filial speech from within. The limit is decisive. Mission and indwelling are not spatial relocation or containment, and an appointed antiphon does not prove which chant a local celebration used.

\Needspace{10\baselineskip}
\noindent\textbf{Can a pardoned people become an inheritance by receiving the One who is given?\enspace\elementref{First Reading, Gospel, Prayer over the Offerings, Communion}}
Exodus ends with Moses asking God to take Israel as the divine inheritance; John centers the Father's gift of the Son; the Communion formula gives the received gift a filial voice; and the offering prayer asks that gifts and worshippers become an enduring offering. The proposed mechanism is a possessive reversal: the people ask to belong to God, God gives rather than grasps, and recipients are enabled to offer themselves. An element-by-element reading can miss how mercy changes the meaning of possession itself. The spiritual fruit is an identity received and freely returned rather than self-owned. The proposal neither treats human offering as the purchase of adoption nor claims that the formulary was composed around a hidden inheritance chiasm.

\Needspace{22\baselineskip}
\noindent\textbf{Does God teach the liturgy how to speak?\enspace\elementref{First Reading, Canticle, Prayer over the Offerings, Communion}}
The Lord first proclaims the divine Name; the furnace canticle answers by blessing that Name; the offering prayer invokes it for sanctification; and the Spirit of the Son finally cries “Abba” in the hearts of children. The mechanism is a transfer of speech agency: self-disclosure generates doxology, invocation, and filial address. What looks like four separate acts of religious language may instead be heard as prayer received before it is uttered. The intellectual and spiritual fruit is a doctrine of worship that resists both verbal mastery and mute passivity: God gives creatures a true answering voice. Yet the four expressions are not lexical equivalents, invocation is not a technique that compels God, and the Spirit's causing filial speech does not dissolve the worshipper's free and ecclesial act.

\Needspace{10\baselineskip}
\noindent\textbf{Can non-condemnation be the condition for serious repair?\enspace\elementref{Second Reading, Gospel, Prayer after Communion}}
John says that the Son is sent to save rather than condemn, while Paul tells a wounded community to repair what is deficient, encourage one another, agree, and live in peace; the final prayer refuses a merely verbal settlement by asking salvation for body and soul. The proposed connection is procedural: persons can face truth and undertake repair because divine action is ordered first toward salvation, not humiliation. An isolated reading can make “not condemn” sound permissive or make correction sound punitive; their conjunction suggests accountable mercy. This may guide reconciliation, safeguarding, and discipline toward restoration without suppressing truth or victims. It cannot erase John's judgment by response to the light, preclude proportionate ecclesial or civil judgment, promise reconciliation without repentance and safety, or turn Paul's closing imperatives into a modern disciplinary code.

\Needspace{10\baselineskip}
\noindent\textbf{Does mercy widen by changing scale rather than changing identity?\enspace\elementref{First Reading, Canticle, Gospel, Prayer after Communion}}
Exodus hears mercy proclaimed within covenant repair; Daniel summons the ordered heights of creation to blessing; John names the world as the object of divine love; and the final prayer asks salvation to reach body and soul. The mechanism is an expanding field---a people, the praising cosmos, the world, the whole human person---held together by the same giver rather than by an abstract theme. Separate treatment can miss that controlled enlargement and the smallness of any Trinitarian devotion indifferent to embodied persons or creation. Its fruit is cosmic praise joined to concrete mercy. But Daniel's canticle arises from a persecuted Jewish narrative, John's “world” cannot simply be equated with a modern ecological category, bodily salvation is not present prosperity, and widening scope neither implies universalism nor dissolves Israel, Church, and creation into one undifferentiated subject.
